\section{Related Work}

\paragraph{Reasoning with LLMs.}

Chain-of-Thought (CoT) \citep{wei-et-al-neurips2022} uses ``reasoning tokens''
to improve the performance of LLMs in many problems.  Many approaches build on
this idea \egcite{yao-et-al-neurips2023,sel2024algorithm}s. For example,
different methods explore multiple reasoning paths to find a valid
solution~\citep{yao-et-al-neurips2023,snell-et-al-icrl2025}. Others show that
using more informative reasoning traces, such as those based on action sequence
similarity, can also enhance planning performance \citep{zhao-et-al-icrl2025}.
However, these approaches often have limited success in generalizing to
out-of-distribution reasoning tasks and longer-horizon planning problems. For
instance, \citet{stechly-et-al-neurips2024} show that while CoT can help with
planning problems, it does not consistently lead to generalization, particularly
for out-of-distribution tasks.

\paragraph{Improving Reasoning Capabilities of LLMs.}

There are several methods in the literature that try to improve the reasoning
capacities of LLMs
\egcite{zhang-et-al-neurips2024,muennighoff-et-al-arxiv2025,wen-et-al-arxiv2025}. The
iterative deployment mechanism we studied here shares conceptual similarities
with the Self-Taught Reasoner (STaR) framework
\cite{zelikman-et-al-neurips2022}, which iteratively improves a model's
reasoning ability by fine-tuning on traces that lead to valid answers.
However, while STaR tries to improve model performance, we focus on a very
different question: we want to understand whether performance improvement could
emerge \emph{unintentionally} from repeated model deployments. We use planning
problems to ground this question in a controlled setting which allows us to
study this phenomenon, and prove that repeated deployment can be seen as a form
of RL. On the technical level, there are key differences between the two as
well.  First, STaR has an extra-step for producing \emph{rationales} for
wrong answers, which are often not available from scraped user-curated
traces. Second, we use valid traces from older generations as well as the
current one emulating the property of web-scraped data containing traces from
multiple former model generations, while STaR focuses on traces from the current
model generation.

Reinforcement learning has been widely used to improve the reasoning
capabilities of large language models
\egcite{shao-et-al-arxiv2024,deepseek-arxiv2025}. Techniques such as group
relative optimization (GRPO) allow us to fine-tune models without supervised
data, using only an internal reward function \cite{shao-et-al-arxiv2024}. By
correctly modeling the reward functions, we can also influence the RL training
so that the models better align with safety conditions or goal specifications.

Iterative deployment is also related to test-time scaling methods. For example,
\citet{muennighoff-et-al-arxiv2025} propose a simple test-time scaling technique
where they fine-tune a small model on traces generated by a much larger, more
capable teacher model---in their paper, the small model is a 32B model, while
the teacher model is a model from the Gemini family. In contrast, iterative deployment
does not rely on a separate, more powerful model. Instead, the LLM generates its
own training data, bootstrapping its performance by solving progressively harder
tasks. The model effectively becomes its own teacher, curating its own
experience to improve in the next generations.


\paragraph{Model Collapse.}
One important connection to our analysis is `model collapse'. \citet{shumailov-et-al-nature2024} show that iteratively training
models with their own synthetic output eventually collapses (i.e., the model's
distribution shrinks so its tails disappear). The key difference in our analysis
is the explicit curation step: only valid traces, as determined by an external
validator, are used for fine-tuning.
\citet{shumailov-et-al-nature2024} lacks a curation step, and assumes all data is kept between generations.
We show that curation can improve reasoning, but it is unknown whether the curation step completely avoids, or simply delays, model collapse. In our context, we have
tested running our experiments above up to 10 generations and (despite the
smaller improvement past the fifth generation) but we did not observe imminent
hints of model collapse within planning domain.


\paragraph{Planning using LLMs.}

\citet{valmeekam-et-al-neurips2023} show that early LLMs do not reliably solve
simple planning tasks. \citet{kokel-et-al-aaai2025} show that problems simpler
than computing a plan, such as verifying if a given plan is valid, are already
challenging for LLMs.

In the context of classical planning, several fine-tuning strategies have been
explored. \citet{rossetti-et-al-icaps2024} train domain-specific GPT models for
each planning domain, achieving strong performance but only for in-distribution
problems. \citet{bohnet-et-al-arxiv2024} report that supervised fine-tuning with
optimal plans produced from an off-the-shelf planner \citep{helmert-jair2006}
did not lead to out-of-distribution
generalization. \citet{huang-et-al-icaps2025} apply RL to end-to-end plan
generation, but observed only a small performance improvement. Our work differs
from these, by showing that generalization can be achieved without access to an
existing planner or to an explicit reward modeling, relying instead on an
iterative process with a simple validation filter.

\citet{verma-et-al-arxiv2025} propose an instruction-tuning framework that
improves symbolic planning by teaching LLMs explicit logical CoT reasoning with
feedback from VAL. While conceptually related to the iterative deployment
mechanism, their method relies on carefully handcrafted training data and a
multi-phase fine-tuning using structured reasoning feedback from a validator,
whereas the mechanism studied here achieves self-improvement through repeated
deployment and fine-tuning solely on valid traces, requiring just a binary
signal.
